\begin{table}[H]
    \centering
    \footnotesize
    \caption{Regressão quantílica linear: ganho na perda \textit{pinball} sobre o quantil incondicional}
    \begin{tabular}{C{2.6cm}C{1.6cm}C{2.6cm}C{1.6cm}C{2.6cm}}
        \toprule
        & \multicolumn{2}{c}{\textbf{Proporção (p.p.)}} & \multicolumn{2}{c}{\textbf{Proficiência (pontos)}} \\
        \cmidrule(lr){2-3}\cmidrule(lr){4-5}
        \textbf{Quantil $\tau$} & \textbf{Ganho (\%)} & \textbf{IC95\% da diferença} & \textbf{Ganho (\%)} & \textbf{IC95\% da diferença} \\
        \midrule
        0,10 & 16,3 & [$-$0,81; $-$0,38] & 20,5 & [$-$1,07; $-$0,52] \\
        0,25 & 13,7 & [$-$1,40; $-$0,48] & 16,0 & [$-$1,79; $-$0,49] \\
        0,50 & 14,6 & [$-$2,03; $-$0,71] & 14,9 & [$-$2,49; $-$0,27] \\
        0,75 & 13,6 & [$-$1,91; $-$0,21] & 10,9 & [$-$1,65; $-$0,04] \\
        0,90 & $-$0,7 & [$-$0,11; 0,18] & 17,0 & [$-$1,39; $-$0,12] \\
        média dos quantis & 12,3 & [$-$1,13; $-$0,44] & 15,1 & [$-$1,52; $-$0,43] \\
        \bottomrule
    \end{tabular}
    \vspace{0.1cm}
    \begin{flushleft}
    {\footnotesize Fonte: elaborada pelo autor, com base nos microdados do Saeb 2023 (Inep). Conjunto de teste. Diferença: perda do modelo menos a do quantil incondicional do treino (negativa: o modelo erra menos). IC95\% por \textit{bootstrap} pareado.}
    \end{flushleft}
    \label{tab:pinball}
\end{table}
